\section{Deployment Considerations}
Although the system is evaluated here as a research prototype, predictive maintenance is ultimately a workflow problem. The implementation includes a lightweight Streamlit dashboard that reads experiment artifacts, visualizes model comparisons, and lets a user inspect benchmark windows through the trained adaptive model. The dashboard returns a health state, maintenance action, coarse subsystem attribution, measured latency, class confidence, and adaptive depth trace for the selected example. That output gives operators a concise summary while still exposing how much internal computation was needed.

\subsection{Operational Flow}
In a live setting, the system ingests sensor batches on a fixed cadence, converts them into sliding windows, and scores each asset independently. If the model repeatedly returns warning or critical states across several windows, the decision policy can escalate to a human operator or a maintenance ticketing system. The early-exit depth can be logged as a diagnostic signal that helps distinguish easy high-confidence cases from ambiguous cases that deserve manual inspection. The dashboard mirrors that flow at demo scale by combining benchmark metrics, model-history plots, confusion matrices, and per-window inference traces in one interface.

\subsection{Reproducibility}
Reproducibility depends on documenting the exact dataset version, preprocessing parameters, random seeds, and train/validation/test split definitions. That matters especially for IMS, where label construction is derived and therefore sensitive to the chosen failure horizon. The released artifact pipeline records the command-line configuration, task-specific metrics, latency benchmarks, and model checkpoints for each variant.

\subsection{Practical Limitations}
There are two main limitations to acknowledge. First, thresholded health labels collapse continuous degradation into coarse categories, which can hide useful structure. Second, the adaptive model may learn to use depth as a confidence proxy in a way that is useful for this benchmark but not necessarily optimal under a different operating regime.
